% Physical page 5 onward: the developed commentary of the concise study,
% organized by the questions the three readings answer differently rather than
% by the order of the Mass. Each unit draws on several appointed elements and
% sets the readings' answers beside one another.

\section{The Propers: Detailed Commentary}
\zlabel{triptych:concise:commentary:start}

\subsection{Which house, and which city?}

The first reading hears the whole Mass from the chants that name a house:
Ps~121:1 at the Introit and again at the Gradual, Ps~101:16 at the Alleluia,
Ps~95:8--9 at the Communion. Its principal witness is St Augustine, who
expounded all three psalms. He begins Ps~121 from its title, a Song of degrees,
``for these degrees are not of descent, but of ascent'', and the ascent is to a
city: ``there is in heaven the eternal Jerusalem, where are our
fellow-citizens, the Angels: we are wanderers on earth from these our
fellow-citizens. We sigh in our pilgrimage; we shall rejoice in the city''
(\work{Enarrationes in Psalmos} 121, 2). His psalter read the first verse
otherwise than the Missal, \latin{Jucundatus sum in his qui dixerunt mihi}, joy
in those who said, and he builds on the persons: \latin{Jucundatus sum in
Prophetis, jucundatus sum in Apostolis}. What his reading gives the Mass is
the plural it shares with the Missal, \latin{ibimus}. No one goes up to this
house alone.

St Hilary of Poitiers reaches the same city by another road. The one who
remembers that he is a future dweller \latin{ciuitatis caelestis uiuis lapidibus
extructae}, of the heavenly city built of living stones, must cry with the
prophet \latin{Laetatus sum}, and Hilary refuses the earthly city in so many
words: \latin{non utique \dots\ hanc terrenam et caducam \dots\ sed illam liberam
et caelestem Hierusalem} (\work{Tractatus super Psalmos}, in Ps.~121, 1--2). He
names its builders, \latin{a Moyse in lege, a prophetis in passionibus suis, a
domino in corpore, ab apostolis in martyriis, a sancto spiritu in uirtutibus},
and gives the plural its reason: \latin{quia unum ecclesiae corpus est}, one
not by a heaping of bodies but \latin{per fidei unitatem, per caritatis
societatem, per operum uoluntatisque concordiam, per sacramenti unum in omnibus
donum} (121, 3, 5). Cassiodorus hears the Holy Spirit as the one who said the
words at which the psalmist rejoices, \latin{qui cordi eius tacita voce
loquebatur}, and the house as \latin{domus de vivis lapidibus fabricata}
(\work{Expositio psalmorum}, in Ps.~121).

The Greek readers of the same psalm read another city. St John Chrysostom sets
Ps~121 at the return from Babylon: the exiles longed for the house of prayer,
and the city built as a city is the city they found in ruins, its towers thrown
down and its walls cast to the ground (\work{Expositio in Psalmum CXXI}, PG 55,
cols.~347--348). Theodoret reads it the same way, and neither reads the Church
or heaven there. St Robert Bellarmine holds both in order, the return from
exile first and then ``the celestial, and not the earthly Jerusalem''; so does
Schuster, commenting on this Mass: ``All this is, of course, to be glorified by
a spiritual interpretation. The peace which is here described is the
atmosphere of the heavenly Jerusalem.'' In the first reading the order is the
same: the return from exile is the Gradual's literal sense and the heavenly
city of Augustine, Hilary and Cassiodorus its allegory. Chrysostom's own turn
to his hearers then serves the literal sense well. Call men to the hippodrome
and many run together; call them to the house of prayer and few are not
sluggish, so that Christians show themselves slower to go up than the returning
Jews were.

The Alleluia brings the nations to the same building. Augustine reads
Ps~101:16 in the psalm's own order: God has had mercy on Sion, his servants have
taken pleasure in her stones, and therefore ``let the heathen fear Thy Name, let
another wall approach also from the heathen, let the Corner Stone be
recognised'' (101, s.~1, 16). The verse the Alleluia stops before gives the work
its present tense: ``This work is going on now. O ye living stones, run to the
work of building, not to ruin'' (101, s.~1, 17). Theodoret, who reads Ps~121 of
the earthly city alone, reads this verse first of the neighbours who marvelled
at Israel's return, and then says that it was fulfilled properly and truly only
after the Incarnation. At this psalm the Greek and the Latin readers meet.

The Communion completes the movement. Augustine reads its psalm under its
title, the house built after the captivity, and at the Communion's verse he
describes the entry: ``Rejoice, because ye have entered into the courts;
rejoice, because ye are being built into the temple of God. For those who enter
are themselves built up, they themselves are the house of God'' (95, 9).
Cassiodorus uses one image, the house of living stones, at all three psalms, and
at this one says that the Church \latin{nec aedificari desinit donec usque ad
finem saeculi praedestinatorum numerus impleatur}, does not cease to be built
until the number of the predestined is filled. In the first reading the request
of the first words, \latin{Da pacem}, is answered in the chants by the
completion of a building.

The Gospel's ``his own city'' enters here too, and it is where the first two
readings meet. The first reads it allegorically with St Hilary: the Gerasene
town that begged Christ to leave \latin{Judaici populi habet speciem}, and,
rejected there, he returns to his own city, for \latin{Deo civitas fidelium plebs
est}, carried in the boat, which is the Church (\work{Commentarius in
Matthaeum} VIII.4). Hilary names no town, and St Thomas Aquinas, in the same
lecture in which he harmonises the Gospels with three cities, Bethlehem by
birth, Nazareth by upbringing and Capernaum by dwelling, reads \latin{in
civitatem gentium, quae sibi datae sunt}. The second reading takes the clause
as the literal place of a verifiable sign, and there the witnesses divide over
the town. Chrysostom says Capernaum, ``that which had Him continually
inhabiting it''; St Jerome says Nazareth, \latin{unde et Nazaræus appellatus
est}; St Augustine, reconciling Matthew with Mark, places the healing at
Capernaum and offers two explanations of the phrase (\work{De consensu
evangelistarum} II.25.58). Nothing in the text settles the town, and nothing
in either reading waits on it: the allegory holds whichever town is meant, and the
literal sign needs only that there was one. The third reading does not use the
clause.

\subsection{What peace is, and what threatens it}

The Gradual's versicle, \latin{Fiat pax in virtute tua}, is the one place where
the Mass says where its peace is to be found, and only the first reading asks
what it means. Augustine asks what a city's strength is and answers with
charity: ``peace be in thy love; for thy strength is thy love'', proving it from
``Love is strong as death'' and drawing the cost: ``this love slayeth what we
have been, that we may be what we were not; love createth a sort of death in
us'' (121, 12). Cassiodorus says the same in almost the same words,
\latin{Virtus quippe ipsius pax sine dubitatione sanctorum est, quae vocatur et
charitas}. Hilary joins peace and strength without making them one:
\latin{pax enim ecclesiae et unitas uirtutem nobis firmitatemque praestat}, the
Church's peace and unity give us strength and firmness (121, 14). On the towers
of the same verse the three give three allegories. For Cassiodorus they are the
martyrs, who defend the city with their bodies; for Hilary the faithful and holy
men who give the rest the firmness of towers; for Augustine the city's
\latin{excelsi}, the few who will sit in judgment and receive the many.

The Epistle shows what threatens such a city, and here two of the readings
draw on the same homily of St John Chrysostom for different things. The Corinthians
are rich ``in all utterance and in all knowledge'', and the verses after the
lesson show them dividing into parties. For the first reading Chrysostom
supplies the remedy. Under the colour of praise Paul ``touches them sharply'',
and ``he always fastens them as with nails to the Name of Christ'', so that
``by incessant application of that glorious Name he may foment their
inflammation, and purge out the corruption of the disease'' (\work{Homilies on
First Corinthians} 2). A city is not held together by its gifts but by the one
Name its members bear. Augustine and Chrysostom agree that the unity is not the
citizens' achievement, and differ in what they are looking at: Augustine
watches a building rise, Chrysostom attends to one coming apart. The Mass sets
the lesson of a divided church between two chants of the city at peace.

For the third reading the same homily supplies its controlling point: the gifts
are given: ``now a favor is not a debt nor a requital nor a payment: which indeed
every where is important to be said, but much more in the case of the
Corinthians who were gaping after the dividers of the Church.'' And on ``the
grace of God'': ``For where grace is, works are not; where works, it is no more
grace. If therefore it be grace, why are you high-minded? Whence is it that you
are puffed up?'' Even the praises, Chrysostom says, only seem to be theirs, ``for
not even did this praise belong to them, but to the grace of God. For that they
had remission of sins, and were justified, this was of the Gift from above''
(2, 1--5). Theodoret and St Thomas Aquinas keep that the gifts are God's and add
that the praise is true. Theodoret says that Paul soothes his hearers before the
cure and yet says nothing false; Aquinas, that the thanks come first \latin{ut
correctionem suorum defectuum tolerabilius ferant}, and that the riches of
utterance and knowledge belong to the more perfect, whom the rest share through
charity (\work{Super I ad Corinthios} c.~1, lect.~1). Ambrosiaster names the
grace given in Christ as the gift by which the believer \latin{gratis accipit
remissionem peccatorum}.

The second reading hears the Epistle differently again, through the continuation
of \work{The Liturgical Year} (Dom Lucien Fromage). The continuation reads the
gifts for which Paul gives thanks, after the September Ember ordinations, as
above all ``the powers conferred by the imposition of the bishop's hands on the
ministers of the Church''. The three uses do not conflict. Chrysostom
himself is the witness both for the one Name that binds and for the grace that
leaves no room for boasting.

\subsection{Power on earth: the unseen proved by the seen}

The second reading's centre is the Gospel's sixth verse, \latin{Ut autem
sciatis, quia Filius hominis habet potestatem in terra dimittendi peccata}. The
forgiveness of the paralytic's sins cannot be verified; his rising from his bed
can, and the second is done so that the first may be believed. Chrysostom reads
the order of the two healings that way. Christ ``healed first that which is
invisible, the soul, by forgiving his sins'', and then offered the visible
miracle as evidence: ``because the one is unseen, the other in sight, I throw in
that, which although an inferior thing, is yet more open to sense; that the
greater also and the unseen may thereby receive its proof'' (\work{Homilies on
Matthew} 29, 1--2). He notices the grammar as well: at first Christ ``said not,
`I forgive thee thy sins,' but, `thy sins be forgiven thee:' upon their
constraining, He discloses His authority more clearly''. St Jerome gives the
argument in fewer words, \latin{Fit igitur carnale signum, ut probetur
spirituale}, and adds that the one who reads the scribes' silent thoughts shows
himself God: \latin{Eadem majestate et potentia qua cogitationes vestras
intueor, possum et hominibus peccata dimittere} (\work{Commentarii in
Matthaeum} I, on 9:4--6).

On why the sins came first, Jerome and Aquinas differ in confidence. Jerome
offers a possibility, \latin{Et idcirco forsan dimittuntur prius peccata tua, ut
causis debilitatis ablatis, sanitas restituatur}: perhaps the sins are forgiven
first so that, the cause of the weakness taken away, health may be restored.
Aquinas states it without the \latin{forsan}: the sins are forgiven first
\latin{quia peccatum erat causa aegritudinis}, as a good physician \latin{qui
causam curat} (\work{Super Matthaeum} c.~9).

St Augustine finds the whole point in one word of address. Matthew's Christ
says ``Son'', Luke's ``Man'', and the difference ``served to indicate that He who
forgave sins to man was Himself God'', while the patient, being a man, could
not say ``I have not sinned'' (\work{De consensu evangelistarum} II.25.57).
Chrysostom and Augustine reach the same conclusion by different roads. Aquinas
then presses the words of verse 6 one by one: \latin{Filius hominis} excludes
Nestorius, \latin{in terra} excludes Photinus, since the power to forgive is
Christ's already on earth. And he raises the objection on which the second
reading turns: the verse seems not to prove Christ's divinity, \latin{quia etiam ipsi
Apostoli habebant potestatem}. His answer is the distinction on which the
second reading rests: \latin{ipsi habebant per viam administrationis, non
auctoritatis}, the Apostles had it by way of ministry, not of authority. He
draws it at the verse that says \latin{in terra}, to keep the authority
Christ's.

The two other readings take the verse otherwise. For the third, what matters is
not the power but its object: a man who could do nothing and was forgiven. For
the first, the Gospel is chiefly the city Christ enters. Only the second hears
the Mass turning on the two words \latin{in terra}, as a statement of where the
power is exercised; a simpler reading of the same texts could remember the
power as Christ's alone, and the difference between them lies in those two
words.

\subsection{Whose faith, and who was carried}

The third reading's Gospel is the man, and on him the witnesses divide in the
open. \latin{Et videns Iesus fidem illorum}: Mt 9:2 says ``their faith'' and no
more. St Jerome reads the phrase in its strongest form: the man was brought
\latin{quia ipse ingredi non valebat}, and \latin{Videns autem Jesus non ejus
fidem qui offerebatur, sed eorum qui offerebant}, Jesus saw not the faith of
the one offered but of those who offered him. He marvels that Christ calls such
a man son, \latin{quem sacerdotes non dignabantur attingere}, and reads the
story of the soul, which, lying powerless, is offered to the Lord by a perfect
teacher and, once healed, \latin{tantum roboris accipit, ut portet statim
lectulum suum} (PL 26, cols.~54--55). St Ambrose, on Luke's telling, takes the
same side, \latin{Magnus Dominus qui aliorum merito ignoscit aliis}, and turns
it to his hearer: \latin{adhibe Ecclesiam quae pro te precetur}, call in the
Church to pray for you (\work{Expositio in Lucam} V.11).

Chrysostom begins from the same point and refuses to stop there. Christ saw
``the faith of them that had let the man down'', but ``in this case the sick man
too had part in the faith; for he would not have suffered himself to be let
down, unless he had believed''; and when the scribes murmur, the man ``gives
himself up to the power of the healer''. Aquinas will not choose: \latin{Curat
aliquando Dominus aliquem propter fidem suam: aliquando propter preces suas, et
aliorum}. The disagreement is real, and the verse does not settle it. What both
sides grant is what the third reading rests on: the man did not come by
himself. Whether the faith Christ saw was the bearers' alone or his as well, it
brought him only as far as he was carried.

St Hilary carries the story over to all humankind. \latin{Jamque in
paralytico gentium universitas offertur medenda}: in the paralytic the whole
body of the nations is offered for healing, brought by angels, called son
because Adam was God's first work, and pardoned the first transgression,
\latin{non enim paralyticum peccasse aliquid accepimus} (\work{Commentarius in
Matthaeum} VIII.5). The nations, too, were brought, and the first sin forgiven
them is Adam's. The third reading hears the Alleluia in the same key: its
nations fear the Lord's name because God has had mercy on Sion and built her,
not because they found him, and Augustine reads the verse in that order, ``Now
that Thou hast pitied Sion \dots\ hence preaching hath increased among the
heathen''.

The Collect says outright what the Epistle implies. \latin{Quia tibi sine te
placere non possumus}: the Latin says more simply than the English ``without thy
help'' that without God himself we cannot please God. Berno of Reichenau
explains the prayer from the Communion that follows it: \latin{Et quia Deo nec
in sacrificio, nec in oblatione hostiarum sine ejus adjutorio placere possumus,
nec in atria ejus introire, nec in aula ejus adorare eum, merito sacerdos ex
sua et nostra voce Deum deprecatur, dicens: Dirigat corda nostra}
(\work{Libellus} V). Schuster calls it a prayer that ``should establish us in
humility. All the good which we do is the work of grace'', and quotes the
question of the same letter as the Epistle, \latin{quid gloriaris, quasi non
acceperis?} (1 Cor 4:7). The Collect has a place in the other two readings as
well: in the first because a city no one can build for himself is still the city
to which the people go up, in the second because the power at work in the Mass is
expressly not ours.

\subsection{``Such power to men'': the one disagreement between two readings}

The Gospel ends with the crowds, who ``feared, and glorified God that gave such
power to men''. Here, and here only, two readings follow different judgements.
St John Chrysostom hears a confession that falls short: Christ ``makes the
invisible evident by that which is in sight. But nevertheless they still creep
upon the earth''. He does not rebuke them but raises their thoughts by his
works: ``Since for the time it was no small thing for Him to be thought greater
than all men, as having come from God'' (\work{Homilies on Matthew} 29, 2). St Ambrose, at Luke's parallel, judges more
severely and in his own way: those who watch the man rise are unbelievers,
\latin{Spectant surgentem increduli}, who \latin{divini operis miracula malunt
timere quam credere}, and \latin{quia non diligebant, calumniabantur}
(\work{Expositio in Lucam} V.15). Chrysostom's crowd glorify God and fall
short; Ambrose's fearers are the calumniators. St Hilary judges otherwise. All
has been concluded in its order, \latin{conclusa sunt omnia suo ordine}, and
the honour is fitting, paid not for what one man could do but because
\latin{potestas hominibus ac via data sit per verbum ejus, et peccatorum
remissionis, et corporum resurrectionis, et reversionis in coelum}, a power
and a way have been given to men through his word (\work{Commentarius in
Matthaeum} VIII.8). The men for whom God is glorified are those who receive the
forgiveness, not those who give it.

The third reading takes Hilary's side, and the choice changes where its Gospel
ends. On Chrysostom's reading the Gospel closes on praise that still falls short
and must be raised. On Hilary's the crowd glorifies God for what men have
received, and the third reading hears in that praise the thanksgiving,
\latin{Gratias ago}, with which the Epistle opens.

Neither Father gives the second reading its own sense of the verse. That
sense is Bl.\ Ildefonso
Schuster's, commenting on this Mass, and it is the third of three. The crowd's
words ``may be understood as a subjective judgement on the part of the people
who had not yet grasped the divine nature of Christ''; they have ``a still deeper
meaning'', the divine nature working the miracle through the human; and
``Further'', the narrative is ``symbolical and prophetic. This power of
remitting sins had to be communicated to men---that is, to the apostles and to
their successors in the priesthood'' (\work{The Sacramentary} III, p.~169). The
continuation of \work{The Liturgical Year} reads the Gospel of the pastors'
prerogative of forgiving sins, names it as the keys exercised in Penance, and
applies the verse's closing words to the Church allegorically: the angels
``sing glory to God, who gave such power to men''.

None of the Fathers read here takes verse 8 of the Church's ministry.
Chrysostom's crowd falls short, Hilary's men are the receivers, Jerome passes
from the risen man to the call of Matthew, and Luke's parallel, on which Ambrose
comments, does not contain the words. The second reading's ministerial sense
stands in four places. Ambrose, glossing the scribes'
question, says that God alone forgives sins, \latin{qui per eos quoque
dimittit, quibus dimittendi tribuit potestatem}, who forgives also through
those to whom he has given the power (V.13); he says it of the scribes'
question, the parallel of Matthew's verse 3, not of the crowd's words. Aquinas
says it at verse 6. Schuster reads it, prophetically, in verse 8. The
continuation places that power in the keys and in Penance, not at the altar.

\subsection{The altar, the exchange and the courts}

At the Offertory the three readings hear one compiled antiphon three ways. For
the first, its point is its last clause, the one that stands verbatim in the
chapter it cites: Moses offers \latin{in conspectu filiorum Israel}, and the
people are present at their own sacrifice. For the third, its point is its
grammar: every finite verb has Moses for its subject, he sanctifies, offers,
immolates and makes the evening sacrifice, and the people are named only as
those in whose sight it is done. For the second, its point is the chapter it
cites, which goes on to the blood sprinkled on altar and people, ``This is the
blood of the covenant'' (Ex 24:8). Chrysostom, expounding that sacrifice as the
Letter to the Hebrews tells it, joins it to the remission of sins: the old rites
were ``not a perfect purification, nor a perfect remission'', but ``in this case
He says, This is the blood of the New Testament, which is shed for you, for the
remission of sins'', and ``With this blood not Moses but Christ sprinkled us''
(\work{Homilies on Hebrews} 16, 3--5). Chrysostom is expounding Hebrews and not
the Offertory's verses; Cornelius a Lapide, on Exodus itself, reads the
covenant blood as the type of the blood Christ named in instituting the
Eucharist. In the second reading the altar meets the Gospel's power through
that blood alone. Schuster regrets that the Missal has cut the antiphon to a single
verse, losing the lines in which the lawgiver ``intercedes for the apostate
people, imploring mercy for them''; the 1962 Missal sings only Moses who
sanctifies and offers.

The Secret divides the same way. For the second reading it claims for a visible
rite an invisible effect, in the present, \latin{efficis}, as the Gospel's sign
claimed an invisible forgiveness; Schuster explains the exchange, ``whilst we
offer him our gifts he, in his turn, bestows on us the gift of himself''. For
the third it divides its verbs: \latin{sicut tuam cognoscimus veritatem} is
indicative and already given, \latin{sic eam dignis moribus assequamur} is
subjunctive and asked for. For the first it names the unity toward which the
city is built, and names it first as God's, \latin{unius summae divinitatis
participes}.

The Communion's \latin{Tollite hostias} names sacrifices, and its witnesses
identify them in two ways that do not exclude each other. The third reading
hears the verse with one of them and the second with the other. Augustine names the one present a creature can bring: ``Confession is a present
unto God \dots\ Enter with an humble heart into the house of God, and thou hast
entered with a present. But if thou art proud, thou enterest empty. For whence
wouldest thou be proud, if thou wert not empty?'' (95, 9). Cassiodorus reads
the \latin{hostiae} as \latin{conscientiae pura libamina, unde non sanguis
currat, sed piae lacrymae defluant}, and notices the order of the imperatives,
\latin{prius posuit, tollite, et sic introite}: first bring, and so enter.
That offering of the heart is the one the third reading hears. Theodoret reads
the rational sacrifices ``we see continually offered and sacrificed by the
priests'', and the expositions under St Athanasius's name the spiritual
sacrifices of the churches' ``priests, I mean, and presidents''; that is the
sacrifice the second reading hears. Neither identification denies the other.
Bellarmine gives both, the contrite heart first and then the Eucharistic
sacrifice offered ``through the hands of the priests of the New Testament'',
and the continuation holds both together: the Communion ``is addressed to the priests, and, at
the same time, to us all'', and with the priest's Victim we bring ``other
victims, that is ourselves''. The first reading hears the same verse as the
entry into courts where ``those who enter are themselves built up''. Schuster
sets the two covenants side by side: ``In the Old Covenant it was the people who
brought gifts to God in his temple, in the New it is God who gives himself to
his people.''

\subsection{Where the Mass sends its people}

The Gospel ends with a man sent home, \latin{vade in domum tuam}, and the
Postcommunion with a prayer, \latin{ut dignos nos eius participatione
perficias}. Each reading ends where the Mass sends its people, and each hears
that prayer asking for something still to come.

The first leaves them at the courts, with the city still ahead. The house the
Introit says we will go into is entered in this Mass in sign, and the psalm's
\latin{ibimus} stays future: the city is the eternal Jerusalem, \latin{unde
nunquam aliquis velit exire} (Cassiodorus), and ``when the building is finished,
and the house dedicated'' the time for running is over (Augustine, on
Ps~101).

The second leaves them forgiven and sent home. Hilary makes the return home the
last step of an order: the remission of sin, the power of resurrection, the end
of the body's weakness and pain, and last the way into paradise, which Adam
left, given back to believers (\work{Commentarius in Matthaeum} VIII.7). Ambrose
calls that home the true one, \latin{quae hominem prima suscepit: non jure
amissa, sed fraude} (V.14). Aquinas sends the forgiven man \latin{in domum
aeternitatis, vel in conscientiam propriam}, and reads the command morally:
\latin{Surge, a peccato per contritionem; tolle lectum, per satisfactionem}.
St Bede, on Luke, reads the bed taken up as the flesh curbed by the reins of
continence (\work{In Lucam} II, on 5:24).

The third leaves them giving thanks for a gift and asking to be made worthy of
it. The Corinthians wait for the Lord who ``will confirm you unto the end
without crime''. The readers divide at that verse: for Chrysostom it is ``also
covertly accusing them'', marking them as still wavering; for Aquinas, for
whom \latin{sine crimine} is \latin{sine peccato mortali}, it is a promise of
perseverance, as it is for Ambrosiaster and Cornelius a Lapide. On either reading the confirming is
the Lord's work. The Postcommunion asks God himself for the fitness, and its
order, the gift, the thanks and then the prayer to become fit for it, is the
third reading in small; Schuster's conclusion is that a fervent Communion ``is
the best preparation for the next''.

Pressed too far, a reading governed by the confession that without God we
cannot please him would become quietism, and the formulary refuses that: the Secret's \latin{dignis moribus assequamur}, the Communion's
\latin{Tollite hostias, et introite} and the Gospel's \latin{Surge, tolle
lectum tuum, et vade} all ask something of the one who prays. Chrysostom, on
the Epistle's ``waiting for the revelation'' (1:7), says that Paul terrifies
them with the judgment seat, ``implying that not only the beginnings must be
good, but the end also''. The Collect answers the objection itself: it says not
that we cannot please God, but that we cannot please him \latin{sine te}. The
man in the Gospel does rise and walk. What he could not do was be forgiven, or
bring himself to where forgiveness was given, and Aquinas notices the shape of
the command: \latin{quia portabatur, praecepit ut portaret}, because he was
carried, he commanded him to carry. The one who was carried now carries what
carried him.

In all three readings the last word of the formulary is the one the
Postcommunion gives. The first hears in it a city still being built, the second
a people forgiven on earth and still to be perfected, the third a gift that
even at the end is asked of God. The Mass that began by asking peace for those
who wait ends by asking to be made worthy of what it has received.
